% ==============================================================================
% T0/FFGFT-Theorie: Dokument 220
% Titel: Casimir-Falsifikationskriterien:
%        Quantitative Vorhersagen im Übergangsbereich d → L_ξ
% Autor: Johann Pascher
% Datum: Mai 2026
% Bezug: Dokumente 091 (Casimir-CMB), 166 (xi-Potenzleiter), 202 (Feldtheorie)
% Anlass: P. Austin (Information Physics Institute), Falsifikationsfrage
% ==============================================================================

\section*{Vorbemerkung}

Dieses Dokument formuliert ein experimentell überprüfbares
Falsifikationskriterium für FFGFT im Casimir-Bereich, in direkter
Antwort auf die methodische Anfrage von P.~Austin (Mail vom 9.~Mai
2026, IPI-Diskussion): \emph{\enquote{Welche Skalenabhängigkeit
oder welcher Residual-Effekt unterscheidet die FFGFT-Vorhersage von
der Standard-$1/d^4$-Form, gegeben dass die gewöhnliche $1/d^4$-Form
asymptotisch reproduziert wird?}}

Die Antwort ist quantitativ und liegt in einem experimentell
zugänglichen Bereich, in dem aktuell keine Präzisionsmessungen
existieren.

% ==============================================================================
\section{Ausgangslage: was FFGFT bestätigt}
% ==============================================================================

FFGFT bestätigt die Standard-Casimir-Theorie im gesamten bisher
gemessenen Bereich. Konkret:

\begin{itemize}[leftmargin=2em,itemsep=2pt,topsep=2pt]
  \item \textbf{Decca et al.\ (170--300\,nm, 0{,}5\,\% Genauigkeit):}
        keine messbare FFGFT-Korrektur erwartet.
  \item \textbf{Bressi et al.\ (0{,}5--3\,$\mu$m):} keine messbare
        FFGFT-Korrektur erwartet.
  \item \textbf{Lamoreaux (0{,}6--6\,$\mu$m):} keine messbare
        FFGFT-Korrektur erwartet.
\end{itemize}

Im Sub-$\mu$m- bis wenige-$\mu$m-Bereich fällt die FFGFT-Korrektur in
die vierte Potenz von $(d/L_\xi)$ ab und ist damit unsichtbar (siehe
§\ref{sec:vorhersage}).

\textbf{Das Falsifikationskriterium liegt im noch unvermessenen
Bereich $d \sim L_\xi \approx 100\,\mu$m}, in dem aktuelle
Drehbalance-Experimente (Iannuzzi-Gruppe) und geplante Messungen
(Dartmouth, ILL Grenoble) gerade beginnen, präzise Daten zu liefern.

% ==============================================================================
\section{Die FFGFT-Übergangsfunktion}
\label{sec:vorhersage}
% ==============================================================================

\subsection{Asymptotische Form (etabliert)}

Aus Dok.~091 ist die Casimir-Energiedichte im klassischen Bereich
$d \ll L_\xi$:
\begin{equation}
  |\rho_\text{Casimir}(d)|
  \;=\; \frac{\pi^2}{240\,\xi_0}\,\rho_\text{CMB}\,
        \left(\frac{L_\xi}{d}\right)^4
  \;=\; \frac{\pi^2 \hbar c}{240\,d^4},
  \label{eq:casimir-asymptotic}
\end{equation}
mit $L_\xi = (\xi_0\hbar c/\rho_\text{CMB})^{1/4} \approx 100{,}24\,\mu$m
und $\xi_0 = 4/30000$.

Das ist die Standardformel — exakt das, was Decca, Bressi, Lamoreaux
gemessen haben, im Rahmen ihrer Genauigkeit.

\subsection{\texorpdfstring{Übergangsfunktion bei $d \to L_\xi$}{Uebergangsfunktion bei d gegen L_xi}}

In FFGFT ist $L_\xi$ keine willkürliche Skala, sondern jene
Vakuum-Längenskala, bei der die Casimir-Energiedichte zwischen
zwei Platten mit der CMB-Energiedichte zusammenfällt:
\begin{equation}
  \rho_\text{Casimir}(d = L_\xi)
  \;=\; \frac{\pi^2}{240}\,\rho_\text{CMB}\,\xi_0^{-1}.
\end{equation}
Bei $d > L_\xi$ kann der Casimir-Mechanismus (Beschneiden der
Vakuummoden zwischen den Platten) nicht mehr operieren, weil die
Wellenlängen der relevanten Moden den Plattenabstand überschreiten
würden.

Daraus folgt eine \textbf{geometrisch motivierte Übergangsfunktion}:
\begin{equation}
  \rho_\text{Casimir}^{\,\text{FFGFT}}(d)
  \;=\; \frac{\pi^2 \hbar c}{240\,d^4}
        \cdot \mathcal{F}\!\left(\frac{d}{L_\xi}\right),
  \label{eq:ffgft-form}
\end{equation}
mit der dimensionslosen Übergangsfunktion
\begin{equation}
  \mathcal{F}(x) \;=\; \frac{1}{1 + x^4},
  \qquad x = d/L_\xi.
  \label{eq:transition}
\end{equation}

\subsection{Begründung der Form $1/(1+x^4)$}

Die Wahl dieser Übergangsfunktion ist nicht willkürlich:
\begin{itemize}[leftmargin=2em,itemsep=2pt,topsep=2pt]
  \item \textbf{Asymptote $x \to 0$ (also $d \ll L_\xi$):}
        $\mathcal{F} \to 1$, Standard-$1/d^4$ wird exakt reproduziert.
  \item \textbf{Asymptote $x \to \infty$ (also $d \gg L_\xi$):}
        $\mathcal{F} \to 1/x^4$, also $\rho_\text{Casimir} \to
        \rho_\text{CMB}$ konstant -- das Vakuum zwischen den Platten
        ist nicht mehr von der CMB unterscheidbar.
  \item \textbf{Übergangspunkt $x = 1$ ($d = L_\xi$):}
        $\mathcal{F}(1) = 1/2$, exakt der Übergang zwischen den
        beiden Regimen.
  \item \textbf{Potenz $4$:} folgt aus der Vierdimensionalität des
        Energiedichte-Volumens $\hbar c/L^4$ und ist konsistent mit
        der CMB-Brückenformel aus Dok.~091.
\end{itemize}

% ==============================================================================
\section{Quantitative Vorhersagen}
\label{sec:zahlen}
% ==============================================================================

Aus (\ref{eq:transition}) folgt die relative Abweichung von der
Standardformel:
\begin{equation}
  \Delta(d) \;\equiv\; \frac{\rho_\text{Casimir}^{\,\text{FFGFT}}(d)
                              - \rho_\text{Casimir}^{\,\text{std}}(d)}
                            {\rho_\text{Casimir}^{\,\text{std}}(d)}
  \;=\; \frac{1}{1 + (d/L_\xi)^4} - 1
  \;=\; -\frac{(d/L_\xi)^4}{1 + (d/L_\xi)^4}.
  \label{eq:delta}
\end{equation}

Mit $L_\xi = 100\,\mu$m ergeben sich folgende Zahlen:

\begin{center}
\renewcommand{\arraystretch}{1.3}
{\small
\begin{tabular}{>{\raggedright\arraybackslash}p{1.6cm}>{\raggedright\arraybackslash}p{1.8cm}>{\raggedright\arraybackslash}p{2.4cm}>{\raggedright\arraybackslash}p{2.4cm}}
\toprule
\textbf{$d$} & $(d/L_\xi)^4$ & \textbf{$|\Delta|$} & \textbf{Messbar?} \\
\midrule
100\,nm     & $10^{-12}$              & $10^{-10}\,\%$    & nein \\
1\,$\mu$m   & $10^{-8}$               & $10^{-6}\,\%$     & nein \\
3\,$\mu$m   & $8{,}1\!\cdot\!10^{-7}$ & $\sim 10^{-4}\,\%$ & nein \\
10\,$\mu$m  & $10^{-4}$               & $\approx 0{,}01\,\%$ & knapp \\
30\,$\mu$m  & $8{,}1\!\cdot\!10^{-3}$ & $\approx 0{,}80\,\%$ & \textbf{ja} \\
50\,$\mu$m  & $0{,}063$               & $\approx 5{,}9\,\%$ & \textbf{ja} \\
70\,$\mu$m  & $0{,}24$                & $\approx 19{,}3\,\%$ & \textbf{ja} \\
100\,$\mu$m & $1{,}0$                 & $50\,\%$          & \textbf{ja} \\
\bottomrule
\end{tabular}
}
\renewcommand{\arraystretch}{1.0}
\end{center}

\subsection{Drei Falsifikationsregime}

\textbf{Regime A: $d < 10\,\mu$m -- Bestätigungsbereich.}
FFGFT sagt $|\Delta| < 0{,}01\,\%$ voraus. Alle bestehenden
Präzisionsmessungen (Decca: 170--300\,nm; Bressi: 0{,}5--3\,$\mu$m;
Lamoreaux: bis 6\,$\mu$m) liegen in diesem Bereich. \emph{FFGFT
widerspricht keinem dieser Experimente.}

\textbf{Regime B: $10\,\mu\text{m} < d < 100\,\mu$m -- Falsifikationsbereich.}
FFGFT sagt eine wachsende, monoton steigende Abweichung von $0{,}01\,\%$
(bei $d=10\,\mu$m) auf $50\,\%$ (bei $d=L_\xi$) voraus. Aktuell
existieren in diesem Bereich keine Präzisionsmessungen, aber die
benötigte Sensitivität ($\sim 1\,\%$ bei $d=30$--$50\,\mu$m) ist mit
Drehbalance-Setups (Iannuzzi-Gruppe, Dartmouth, ILL Grenoble)
erreichbar.

\textbf{Regime C: $d \to L_\xi$ -- Vorhersagepunkt.}
Bei genau $d = L_\xi$ sollte $\rho_\text{Casimir} = (1/2)\rho_\text{Casimir}^\text{std}$
gemessen werden, mit $L_\xi$ aus Dok.~091 zu $100{,}24\,\mu$m
festgelegt. Dies ist der \emph{schärfste} Test: ein einziger
Messpunkt bei $100\,\mu$m mit $\sim 5\,\%$ Genauigkeit fixiert oder
widerlegt die FFGFT-Vorhersage.

% ==============================================================================
\section{Konkrete Falsifikationskriterien}
% ==============================================================================

FFGFT wird durch die folgenden Beobachtungen \textbf{widerlegt}:

\begin{enumerate}[leftmargin=2em,itemsep=4pt,topsep=2pt]
  \item \textbf{Sub-$\mu$m-Bereich:}
        Eine Casimir-Messung bei $d < 1\,\mu$m, die eine Abweichung
        $|\Delta| > 0{,}1\,\%$ vom Standard-$1/d^4$ zeigt.
        \emph{Begründung:} FFGFT sagt hier $|\Delta| < 10^{-6}\,\%$
        voraus.

  \item \textbf{Übergangsbereich:}
        Eine Präzisionsmessung im Bereich $30$--$70\,\mu$m, die
        \emph{keine} systematische negative Abweichung in der
        durch (\ref{eq:delta}) vorhergesagten Größe zeigt.
        \emph{Begründung:} FFGFT sagt hier $|\Delta|$ zwischen
        $1\,\%$ und $20\,\%$ voraus, mit eindeutigem Vorzeichen
        (Casimir-Kraft schwächer als $1/d^4$).

  \item \textbf{Vorhersagepunkt:}
        Eine Messung bei $d = 100\,\mu$m, die nicht
        $\rho_\text{Casimir}(L_\xi) \approx \frac{1}{2}\,
        \rho_\text{Casimir}^\text{std}(L_\xi)$ liefert
        (innerhalb der Messgenauigkeit), oder
        $L_\xi$ deutlich verschieden von dem aus
        $\rho_\text{CMB}$ berechneten Wert
        $L_\xi^\text{T0} = 100{,}24\,\mu$m liefert.
        \emph{Begründung:} Die Brücke zur CMB
        (Dok.~091, Gl.~\ref{eq:cmb_vacuum_relation})
        wäre damit gebrochen.

  \item \textbf{Funktionale Form:}
        Eine systematische Abweichung der gemessenen
        Übergangsfunktion von der Form
        $\mathcal{F}(d/L_\xi) = 1/(1 + (d/L_\xi)^4)$
        außerhalb der Messunsicherheit.
\end{enumerate}

Jedes dieser Kriterien ist hinreichend für die Widerlegung der
FFGFT-Casimir-Vorhersage. Sie sind voneinander \emph{unabhängig}, so
dass mehrere Experimente parallel oder sequenziell die Theorie
testen können.

% ==============================================================================
\section{Was FFGFT NICHT vorhersagt}
% ==============================================================================

Zur methodischen Klarheit:

\begin{itemize}[leftmargin=2em,itemsep=2pt,topsep=2pt]
  \item FFGFT sagt \emph{keine} Yukawa-artige Zusatzkraft
        $\sim e^{-d/\lambda}$ voraus. Modifikationen aus Extra-Dimensionen
        (Randall-Sundrum, ADD) liefern solche Korrekturen; FFGFT
        nicht. Ein gemessener Yukawa-Term ist kein Beweis und keine
        Widerlegung von FFGFT.
  \item FFGFT sagt \emph{keine} Modifikation der thermischen
        Casimir-Korrektur (Drude vs.\ Plasma-Modell-Streit) voraus.
        Die FFGFT-Vorhersage steht orthogonal zu dieser Debatte
        und ist mit beiden Modellen kompatibel.
  \item FFGFT sagt \emph{keine} Modifikation der Geometriefaktoren
        (Sphäre-Platte vs.\ Platte-Platte, Proximity Force
        Approximation) voraus. Die Korrektur (\ref{eq:transition})
        ist universell in $d/L_\xi$ und sollte in allen Geometrien
        gleich auftreten.
\end{itemize}

% ==============================================================================
\section{Zusammenfassung}
% ==============================================================================

FFGFT macht im Casimir-Bereich eine quantitative, experimentell
prüfbare Vorhersage:
\begin{equation*}
  \boxed{
    \rho_\text{Casimir}^{\,\text{FFGFT}}(d)
    \;=\; \rho_\text{Casimir}^{\,\text{std}}(d)
          \cdot \frac{1}{1 + (d/L_\xi)^4},
    \qquad L_\xi = 100{,}24\,\mu\text{m}.
  }
\end{equation*}

Die Vorhersage:
\begin{itemize}[leftmargin=2em,itemsep=2pt,topsep=2pt]
  \item ist im bestehenden Messbereich ($d < 10\,\mu$m) mit allen
        Daten kompatibel (Bestätigungsregime, $|\Delta| < 0{,}01\,\%$);
  \item liefert im Bereich $30$--$100\,\mu$m eine messbare,
        monoton steigende Abweichung ($0{,}8$--$50\,\%$);
  \item ist mit aktueller Drehbalance-Technologie experimentell
        zugänglich;
  \item enthält keinen freien Parameter -- $L_\xi$ ist über
        Dok.~091 numerisch durch $\rho_\text{CMB}$ festgelegt;
  \item ist bei vier unabhängigen Kriterien falsifizierbar
        (Sub-$\mu$m, Übergangsbereich, Vorhersagepunkt,
        funktionale Form).
\end{itemize}

\textbf{Methodische Einordnung.} Diese Vorhersage ist
FFGFT-intern. Sie folgt aus der $\xi$-Vakuum-Geometrie und der
Casimir-CMB-Brückenformel (Dok.~091), nicht aus einer
topologischen Partition oder einer rekursiven Survival-Struktur.
Die Frage, ob andere IPI-Stränge dieselbe Vorhersage in ihrer
Sprache reproduzieren können, bleibt offen und ist nicht
Voraussetzung für die Gültigkeit der Aussage.

\textbf{Status.} Die Vorhersage steht zur experimentellen Prüfung
zur Verfügung. Falls bestehende oder kommende Drehbalance-Daten
im $30$--$100\,\mu$m-Bereich verfügbar werden, sind sie direkt
mit (\ref{eq:transition}) zu vergleichen.
